\section{Method Overview}
\label{sec:method-overview}

We model the cognitive state of a French speaker as a tuple
$\rho = (\rho_{\mathrm{phono}}, \rho_{\mathrm{sem}},
\rho_{\mathrm{lex}}, \rho_{\mathrm{syntax}}) \in (\Delta^{32})^4$
--- four discrete probability simplices, one per Levelt-Baddeley
species, each ranging over 32 activation stacks. A French query
$q$ is encoded by a text encoder $f_\theta$ into a 384-dimensional
embedding. The objective is a free energy $F(\rho, q)$ whose
gradient drives a JKO step that produces a state $\rho'$
reflecting how the query reshapes activation across species.

\paragraph{Pipeline.} The coupling proceeds in five stages
(\cref{fig:pipeline}). (1)~A query $q$ is passed through
$f_\theta$, producing a 384-d embedding
$q_{\mathrm{emb}} \in \mathbb{R}^{384}$. (2)~The current state
$\rho_\tau \in (\Delta^{32})^4$ is flattened to a 128-d vector
$\rho_\tau^{\mathrm{flat}} \in \Delta^{128}$ and decoded by the
JEPA predictor $g_{\mathrm{JEPA}}$ to produce a reconstruction
$\hat{q} = g_{\mathrm{JEPA}}(\rho_\tau^{\mathrm{flat}})$.
(3)~The accuracy term
$\|q_{\mathrm{emb}} - \hat{q}\|^2$ is differentiated via JAX
autodiff; the complexity and coupling terms are analytical.
(4)~The gradient $\nabla F(\rho_\tau, q)$ is combined with the
Wasserstein proximal operator to produce one JKO step (see
Section~\ref{sec:query-conditioned-f}). (5)~The updated state
$\rho_{\tau+1}$ is projected back onto the simplex via
\texttt{softplus}-normalize.

\begin{figure}[t]
  \centering
  \small
  \begin{tabular}{c}
    \fbox{\parbox{0.9\linewidth}{\centering
      \texttt{query $q$} \quad $\Longrightarrow$ \quad
      $f_\theta$ \quad $\Longrightarrow$ \quad
      $q_{\mathrm{emb}} \in \mathbb{R}^{384}$ \\[1ex]
      $\rho_\tau \in (\Delta^{32})^4$ \quad $\Longrightarrow$ \quad
      \texttt{flatten} \quad $\Longrightarrow$ \quad
      $g_{\mathrm{JEPA}}$ \quad $\Longrightarrow$ \quad $\hat{q}$ \\[1ex]
      \bigskip
      $\nabla F(\rho_\tau, q) \;=\;
       \nabla D_{\mathrm{KL}} \;+\;
       \nabla \tfrac{1}{2\sigma^2}\|q_{\mathrm{emb}} - \hat{q}\|^2 \;+\;
       \nabla \lambda_J \sum J_{st}\langle\rho_s, \rho_t\rangle$ \\[1ex]
      $\Downarrow$ \quad \texttt{JKO step + simplex-project} \\[1ex]
      $\rho_{\tau+1} \in (\Delta^{32})^4$
    }}
  \end{tabular}
  \caption{\textsc{QueryConditionedF} pipeline. Analytical
    gradients (KL, species coupling) compose with autodiff
    through the JEPA accuracy term; the JKO step projects the
    update back onto the product simplex.}
  \label{fig:pipeline}
\end{figure}

\paragraph{Design choices.} Three are load-bearing. First,
\emph{32 bins per species} is a compromise between representational
resolution (32 preserves enough granularity for the lemma-hash
semantic mapping of \S\ref{sec:labeler}) and numerical stability
of the JKO step (gradients through very-high-$K$ simplices require
log-space normalization with attendant noise floors). Second,
\emph{four species} follow the Levelt-Baddeley decomposition
\citep{levelt1989speaking, baddeley1974} verbatim; our functional
does not rely on the specific count --- any finite species number
with pre-defined priors would fit identically. Third, \emph{JKO
over Euler} is motivated by the Wasserstein geometry of the
simplex: the squared Euclidean gradient step would systematically
shrink entropy and collapse onto the prior, whereas JKO respects
the geometry of the underlying probability measure.

Section~\ref{sec:query-conditioned-f} defines $F$ and its
gradients formally. Section~\ref{sec:experiments} reports the
empirical ablation: three encoder architectures train against
$\rho \to \rho' - \rho$ pairs produced by the JKO oracle on a
50\,k hybrid French corpus.
